% =====================================================================
%  Section VI. Phase Transitions
% =====================================================================
\section{PHASE TRANSITIONS}\label{sec:phases}

In the strong-coupling limit ($g\to\infty$, $\kappa\to0$) the gauge
theory is far from being Lorentz-invariant.  More precisely, since the
action was defined on a Euclidean metric, it is Euclidean invariance that
is missing.  In the strong-coupling limit, vacuum expectation values
decrease rapidly at separations of only a few lattice sites (there is a
factor $g^{-2}$ or $\kappa$ or both for each unit lattice spacing of
separation).  This corresponds to the existence of masses much larger
than the cutoff.  [The usual rule is that if a propagator falls as
$e^{-\xi x}$ for $x$ large then the lowest mass intermediate state
contributing to the propagator has mass $1/\xi$.  If the propagator
behaves as $g^{-2x/a}$ for distances $x=na$, then the corresponding mass
is $2(\ln g)/a$.  This is larger than the cutoff momentum $\pi/a$ if $g$
is large.]

Thus, one is interested in practice in values of $g$ and $\kappa$ such
that the correlation length $\xi$ is much larger than the lattice
spacing $a$, in order that the corresponding mass is much less than the
cutoff.  One knows from statistical mechanics that large correlation
lengths are associated with second-order phase transitions (critical
points).  Thus one seeks special values $g_{c}$ and $\kappa_{c}$ for $g$
and $\kappa$ at which there is a phase transition.

It has already been argued that there are two distinct phases for the
gauge field, a strong-coupling phase for large $g$ which binds quarks,
and a weak-coupling phase for small $g$ which does not bind quarks.  The
arguments given neglected quark vacuum loops, which is reasonable if
$\kappa$ is small.  There should be a transition between these two
phases which would occur at a critical value $g_{c}$ for any $g$ and any
small value of $\kappa$.  This is one possible phase transition; it is
this transition which was discussed in Sec.~\ref{sec:binding}.  But, as
will be argued below, this is probably a first-order transition rather
than second order.

Suppose one wishes to construct a model of strong interactions using the
lattice theory of this paper with the gauge group separate from ordinary
SU(3)$\times$SU(3) symmetry.  Then the gauge fields would all be
SU(3)$\times$SU(3) singlets, while the quark fields would carry
SU(3)$\times$SU(3) quantum numbers as well as gauge-group quantum
numbers.  A little thought shows that in the strong-coupling limit
($g\to\infty$, $\kappa\to0$), SU(3)$\times$SU(3) is an exact symmetry
rather than a spontaneously broken symmetry.  Varying $g$ does not change
this situation; so one must hope that by increasing $\kappa$ one can
change the exact SU(3)$\times$SU(3) into broken SU(3)$\times$SU(3).  If
this does not work one is free to introduce additional terms into the
quark field action in hopes of forcing a spontaneous breaking of
SU(3)$\times$SU(3).  Suppose, for simplicity, that SU(3)$\times$SU(3)
can be broken by increasing $\kappa$.  Then there will be a phase
transition at a critical value $\kappa_{c}$ for $\kappa$ where one changes
from exact SU(3)$\times$SU(3) to spontaneously broken
SU(3)$\times$SU(3).  If this transition is a second-order transition then
there will be a large correlation length for $\kappa$ near $\kappa_{c}$;
in this case the theory might be a realistic model of broken
SU(3)$\times$SU(3) for $\kappa$ slightly greater than $\kappa_{c}$ (with
$g$ large enough to maintain quark binding).

In summary, the transition of real interest is a transition in $\kappa$
(or some other parameter introduced into the quark action) rather than
$g$.

Apart from special limits ($g\to\infty$ and $\kappa\to0$, or $g$ small)
it is very difficult to solve the lattice theory.  It is especially
difficult to solve the lattice theory near a critical point with a large
correlation length.  Various methods have been developed by statistical
mechanicians to deal with this problem.  In the remainder of this section
these methods will be discussed briefly.  There are essentially three
approaches to consider:  (1) mean-field techniques, (2) series
expansions, and (3) the renormalization-group approach.

Mean-field techniques${}^{15}$ are the simplest and crudest methods for
studying a critical point; invariably they are the methods one uses
first in studying a new situation.  They are used to determine if there
is a phase transition, whether it is first or second order, and to give
rough estimates of the behavior near the transition.  None of the
results of a mean-field calculation are entirely trustworthy.  Examples
of mean-field calculations will be given later.

An example of a series expansion would be the expansion of the
current-current propagator for small momentum (momentum $\ll 1/a$) in
powers of $g$ and $\kappa$, to high order in $g^{-2}$ and $\kappa$.  One
then uses Pad\'e-approximant techniques to look for singularities in
either $g$ or $\kappa$ that would be associated with a mass approaching
0.  In simple statistical-mechanical problems one can generate 12 terms
or so in analogous expansions.  The expansion for the lattice theory of
this paper is more complicated, but one could hope to generate maybe 6
or 7 orders with some practice.  Series expansions require considerably
more effort than mean-field calculations; they apply mainly to
propagators, being very awkward to perform on three- and four-point
functions, and one must have a clear idea of what one is trying to learn
before attempting such calculations.  See Ref.~16 for one of the best
series-expansion formalisms; see Ref.~13 for a general review.

The renormalization-group approach is potentially the most powerful and
accurate method for studying lattice theories near a critical point, but
at present the renormalization-group techniques are too limited in scope
to be applicable to the present problem.  See Refs.~6 and~17.

Return to mean-field ideas.  The prototype mean-field calculation is a
calculation of the magnetization as a function of the external field for
an Ising ferromagnet.  Let $s_{n}$ be the spin at site $n$ with values
$\pm 1$ only; let the interaction be
\begin{equation}
A = K\sum_{n,\mu}s_{n}s_{n+\hat\mu} + h\sum_{n}s_{n},
\label{eq:6.1}
\end{equation}
where $K$ is related to the spin-spin coupling and $h$ is proportional
to the external field.  Then
\begin{equation}
M = Z^{-1}\Bigl(s_{0}\exp\Bigl[K\sum_{n,\mu}s_{n}s_{n+\hat\mu}
    + h\sum_{n}s_{n}\Bigr]\Bigr),
\label{eq:6.2}
\end{equation}
where $(\ )$ means a sum over all configurations of all spins, and
\begin{equation}
Z = \sum\exp\Bigl[K\sum_{n,\mu}s_{n}s_{n+\hat\mu}
    + h\sum_{n}s_{n}\Bigr].
\label{eq:6.3}
\end{equation}

In the mean-field approximation, one assumes that the spins
$s_{n,\hat\mu}$, coupled to $s_{0}$, can be replaced by their average
value $M$.  As a result, the formula for $M$ simplifies to a sum over
$s_{0}$ only, namely
\begin{equation}
M = Z_{0}^{-1}\sum_{s_{0}=\pm 1}s_{0}\exp\bigl[(2dKM+h)s_{0}\bigr],
\label{eq:6.4}
\end{equation}
with $d$ being the dimensionality (3 usually) and
\begin{equation}
Z_{0} = \sum_{s_{0}=\pm 1}\exp\bigl[(2dKM+h)s_{0}\bigr].
\label{eq:6.5}
\end{equation}
The result is
\begin{equation}
M = \tanh(2dKM+h).
\label{eq:6.6}
\end{equation}
If $2dK<1$ this equation has a unique solution for $M$ as a function of
$h$; in particular, $M=0$ for $h=0$.  For $2dK>1$ the solution is
multiple-valued; stability considerations show that one must choose a
solution with $M\ne0$ when $h=0$.

In this approximation one has actually replaced $\sum_{\mu}s_{n,\hat\mu}$
by $2dM$, which is a good approximation if $d$ is large.  This is
generally true of mean-field theories.

An analogous mean-field calculation can be performed for the lattice
gauge theory.  In this case a simple external field term has the form
\[
\frac{1}{g}\sum_{n,\mu}\bigl(e^{i\thl{0\mu}(n)}+e^{-i\thl{0\mu}(n)}\bigr)
\]
(this is to be added to the gauge-field action), and $M$ can be defined
to be the expectation value of $e^{i\thl{0\mu}}$.  The question is
whether $M$ is zero in the limit $h\to0$.  In the gauge-field theory
$e^{i\thl{0\mu}}$ couples to a product of three other exponentials; as a
mean-field approximation one replaces this product by $M^{3}$.  The
result of this is that
\begin{equation}
M = Z_{\theta}^{-1}\int\ddth{0\mu}\,e^{i\thl{0\mu}}
  \exp\bigl[2(d-1)M^{3}+h\bigl(e^{i\thl{0\mu}}+e^{-i\thl{0\mu}}\bigr)\bigr],
\label{eq:6.7}
\end{equation}
where $d$ is the space-time dimensionality,
\begin{equation}
Z_{\theta} = \int\ddth{0\mu}
  \exp\bigl[2(d-1)M^{3}+h\bigl(e^{i\thl{0\mu}}+e^{-i\thl{0\mu}}\bigr)\bigr],
\label{eq:6.8}
\end{equation}
The result of this is that
\begin{equation}
M = f\bigl[2(d-1)M^{3}+h\bigr],
\label{eq:6.9}
\end{equation}
where $f$ is a ratio of Bessel's functions.  If $g$ is large the
solution to this equation is unique and $M=0$ for $h=0$.  If $g$ is
small then there are solutions with $M\ne0$ for $h=0$, and stability
considerations show again that the $M\ne0$ solutions are preferred.

In the magnetic case, one finds that the spontaneous magnetization $M$
goes to zero for $2dK-1\to0$ [from Eq.~\eqref{eq:6.6}].  However, the
gauge-field case never has a solution for $h=0$ with $M$ small but
nonzero.  Thus there is a first-order transition at the value of $g$ for
which $M$ changes from zero to being nonzero.

A nonzero value of $M$ in the limit $h\to0$ means one has spontaneous
breaking of the gauge-field symmetry.  So for small $g$ the theory shows
spontaneous breaking.

A much more thorough discussion of the mean-field approximation has been
given by Balian, Drouffe, and Itzykson.${}^{17}$  A Hamiltonian
formulation of the lattice gauge theory has been given by Kogut and
Susskind.${}^{18}$  A clear review of quark confinement in the lattice
theory is given in Ref.~20.  Another formulation of the connection
between strongly coupled gauge theories and string models is given in
Ref.~21.
